% Einheit 3 von 3 – Argumente über Formel und Verträglichkeit: die Länge des Intervalls bei doppeltem Umfang (Faktor eins durch Wurzel zwei (bank/konfidenzintervalle/e3.jsonl, zusammenbau v0.1)
%% TODO Zeitmarke Einheit 3: Marken-Zeile nicht lesbar
%% TODO Zweigzeile Teil 1: was hier gelernt wird (Satz in Schülersprache aus dem Einheitstitel) – Einheit 3
\einheitenkopf[e3]{Einheit 3 von 3 · Argumente über Formel und Verträglichkeit: die Länge des Intervalls bei doppeltem Umfang (Faktor eins durch Wurzel zwei}
\zweigzeile{keine Prüfungsaufgabe}
\setcounter{aufgabe}{13}

%% TODO Titel: Ich-kann-Satz fehlt in der Bank (Platzhalter: Kettenname) – Nr. 14
%% TODO Anweisung: ein Satz über der ersten Teilaufgabe fehlt in der Bank – Nr. 14
\begin{aufgabe}{Argumente}
%% TODO \rechenplatz{n}: Schrittzahl des Grundfalls fehlt in der Bank (nur bei fester Schreibform, 2.3 b)
\begin{teile}
\teil Stimmt die Aussage? Nur ankreuzen: „Bei vierfachem Umfang und gleichem Stichprobenanteil wird das Konfidenzintervall etwa halb so lang.“ \\ \kreuz{stimmt} \\ \kreuz{stimmt nicht}
\teil Der Stichprobenanteil ist $0{,}35$, die Annahme $p = 0{,}3$. Begründe ohne Rechnung, dass die obere Grenze des Konfidenzintervalls größer als die Annahme ist.
\teil Ein Konfidenzintervall ist bei Umfang 400 etwa $0{,}08$ lang. Wie lang ist es etwa bei gleichem Stichprobenanteil und Umfang 1600? etwa \leerfeld
\teil Der Stichprobenanteil ist $0{,}32$, die Sicherheitswahrscheinlichkeit 95\,\%. Aussage: „Obwohl $0{,}32$ genau in der Mitte zwischen $0{,}28$ und $0{,}36$ liegt, kann $0{,}28$ unverträglich und $0{,}36$ verträglich sein.“ Rechnung I: $0{,}32 = 0{,}28 + 1{,}96 \cdot \sqrt{0{,}28 \cdot 0{,}72 / n}$ ergibt $n \approx 484{,}0$; Rechnung II: $0{,}32 = 0{,}36 - 1{,}96 \cdot \sqrt{0{,}36 \cdot 0{,}64 / n}$ ergibt $n \approx 553{,}2$. Beurteile die Aussage mit beiden Rechnungen \hfill (Abitur 2025 LK).
\end{teile}
\end{aufgabe}

%% TODO Titel: Ich-kann-Satz fehlt in der Bank (Platzhalter: Kettenname) – Nr. 15
%% TODO Anweisung: ein Satz über der ersten Teilaufgabe fehlt in der Bank – Nr. 15
\begin{aufgabe}{Argumente – Fehler finden, Begründen}
\begin{teile}
\teil Bei Umfang 500 ist ein Konfidenzintervall $0{,}1$ lang. Tim sagt: „Bei Umfang 1000 und gleichem Anteil ist es $0{,}05$ lang.“ Finde den Fehler und gib die Länge richtig an.
\teil Begründe, warum der Stichprobenanteil immer im Konfidenzintervall liegt.
\end{teile}
\end{aufgabe}
